\section{Homomorphisms, Ideals and Quotients}

\begin{definition}[Ring Homomorphism]
    \label{def:ring-homomorphism}%
    Let \(R\) and \(S\) be rings. A function \(\functiondomain{f}{R}{S}\) is a \emph{ring homomorphism} if
    \begin{enumerate}
        \item \(f\pqty{r_1 + r_2} = f\pqty{r_1} + f\pqty{r_2}\);
        \item \(f\pqty{0} = 0\);
        \item \(f\pqty{r_1 r_2} = f\pqty{r_1} f\pqty{r_2}\);
        \item \(f\pqty{1} = 1\).
    \end{enumerate}
\end{definition}

\begin{remark}
    Point 4 is \textbf{not implied by} Point 3 (unlike Point 2 is implied by Point 1) since addition and multiplication are not symmetric -- addition gives a group structure, but multiplication does not (we cannot invert multiplication).
\end{remark}

\begin{definition}[Ring Isomorphism, Kernel, Image]
    \label{def:ring-isomorphism-kernel-image}%
    If \(\functiondomain{f}{R}{S}\) is a bijective ring homomorphism, it is an \emph{isomorphism}.

    The \emph{kernel} and \emph{image} of \(f\) is the kernel and image of the corresponding group homomorphism (when the operation is restricting to addition), i.e.,
    \[
        \ker f = \set{r \in R}{f \pqty{r} = 0}
    \]
    and
    \[
        \img f = \set{s \in S}{\exists r \in R: s = f\pqty{r}}.
    \]
\end{definition}

\begin{remark}
    The inverse of a ring isomorphism is also a ring isomorphism.
\end{remark}

\begin{lemma}[Homomorphisms are Injective if and only if it has Trivial Kernel]
    \label{lem:homomorphisms-injective-trivial-kernel}%
    A ring homomorphism \(\functiondomain{f}{R}{S}\) is injective if and only if \(\ker f = \Bqty{0}\).
\end{lemma}

\begin{proof}
    This is implied from the statement on group homomorphisms.
\end{proof}

Recall a group homomorphism's kernel must be a \emph{normal subgroup} -- there is a very similar key notion in rings.

\begin{definition}[Ideal]
    \label{def:ideal}%
    A subset \(I \subseteq R\) is called an \emph{ideal}, written as \(I \idealsub R\), if
    \begin{enumerate}
        \item \(I\) is a subgroup of \(R\) under addition; and
        \item \emph{Absorbing Rule.} If \(a \in I\) and \(b \in R\), then \(ab \in I\).
    \end{enumerate}
\end{definition}

\begin{remark}
    An ideal is not a subring, since it not necessarily contains \(1 \in R\). Indeed, if \(1 \in I \idealsub R\), then \(R = I\).
\end{remark}

\begin{lemma}[Kernels are Ideals]
    \label{lem:kernel-are-ideal}%
    If \(\functiondomain{\phi}{R}{S}\) is a homomorphism, then \(\ker \phi \idealsub R\) is an ideal.
\end{lemma}

\begin{proof}
    Since \(\phi\) is a group homomorphism, \(\ker \phi\) is a subgroup under addition.

    If \(a \in \ker \phi\) and \(b \in R\), then
    \[
        \phi\pqty{ab} = \phi\pqty{a} \cdot \phi\pqty{b} = 0 \cdot \phi\pqty{b} = 0
    \]
    so \(ab \in \ker \phi\), which finishes the proof.
\end{proof}

\begin{remark}
    Let \(u \in R\) is a unit, and \(I \idealsub R\) is an ideal. If \(u \in I\), then \(I = R\). This is since \(1 = u v\) for \(u \in I\) and \(v \in R\), then \(1 \in I\).
\end{remark}

\begin{example}[Ideals of \(\ZZ\)]
    \label{ex:ideals-of-z}%
    Let \(\ZZ\) be the ring of integers. Then
    \[
        n \ZZ = \Bqty{\ldots, -2n, -n, 0, n, 2n, \ldots}
    \]
    is an ideal. In fact, every (non-zero) ideal of \(\ZZ\) has this form, for some \(n \in \ZZ\). This result follows from the subgroup structure of \(\ZZ\), but we reproduce the proof of follows to motivate an algebraic structure.

    Indeed, if \(I \neq \Bqty{0}\), let \(n \in \ZZ\) be the smallest positive element \(n\) that lies in \(I\). We claim \(I = n\ZZ\).

    By the Euclidean algorithm, \(m = q \cdot n + r\) for \(q \in \ZZ\) and \(0 \leq r < n\), but \(q \cdot n \in I\) by definition of an ideal, so \(r \in I\). If \(r \neq 0\), then this contradicts minimality. So \(r = 0\) and this finishes the proof.
\end{example}

But we may generalise this idea -- this only required the Euclidean algorithm -- and we may abstract this to algebraic structures where such an `Euclidean algorithm' exists.

\begin{definition}[Ideal Generated by a Set, Principal Ideal]
    \label{def:ideal-generated-principal}%
    Let \(A \subseteq R\) be a subset. The \emph{ideal generated by \(A\)} is given by
    \[
        \pqty{A} = \set{\sum_{a \in A} r_a a}{r_a \in R, r_a = 0 \text{ for all but finitely many }a}.
    \]

    An ideal \(I \idealsub X\) is called \emph{principal} if \(A = \pqty{r}\) for some \(r \in R\).
\end{definition}

For example, as we see above, all ideals of \(\ZZ\) are principal, since \(n\ZZ\) is generated by \(n\).

\begin{example}[Ideal Generated]
    \label{ex:ideal-generated}%
    The subset of \(\RR\bqty{x}\) given by
    \[
        \set{f\pqty{x} \in \RR\bqty{x}}{f\pqty{0} = 0},
    \]
    i.e. polynomials with zero constant term is an ideal. It is generated by the element \(x \in \RR\bqty{x}\).
\end{example}

\begin{definition}[Quotient Ring]
    \label{def:quotient-ring}%
    Let \(I \idealsub R\). The \emph{quotient ring} is the set of cosets of \(I \in R\), taking the form \(\pqty{r + I}\) with \(0\) given by \(0_R + I\) and \(1\) given by \(1_R + I\), and operations
    \begin{align*}
        \pqty{r_1 + I} + \pqty{r_2 + I} &= \pqty{r_1 + r_2} + I, \\
        \pqty{r_1 + I} \cdot \pqty{r_2 + I} &= \pqty{r_1 \cdot r_2} + I.
    \end{align*}
\end{definition}

\begin{proposition}[Quotient Ring is a Ring]
    \label{prop:quotient-ring-ring}%
    The quotient ring \(\flatfrac{R}{I}\) is a ring, and the function
    \[
        \function{q_I}{R}{\flatfrac{R}{I}}{r}{r + I}
    \]
    is a ring homomorphism.
\end{proposition}

\begin{proof}
    Obviously, \(q_I\) is a well-defined group homomorphism and addition is well-defined. We verify that multiplication is well-defined (since it would suggest some motivation into why an ideal is defined in this particular way). Let
    \[
        r_1 + I = r_1' + I \qand r_2 + I = r_2' + I.
    \]

    Then \(r_1 - r_1' \in I\), say \(r_1 - r_1' = a_1\), and similarly let \(r_2 - r_2' = a_2 \in I\).

    Then
    \[
        r_1' r_2' = \pqty{r_1 + a_1} \pqty{r_2 + a_2} = r_1 r_2 + r_1 a_2 + r_2 a_1 + a_1 a_2.
    \]

    But \(r_1 a_2, r_2 a_1, a_1 a_2 \in I\) by absorption property. Therefore, we can conclude that
    \[
        r_1 r_2 + I = r_1' r_2' + I,
    \]
    so multiplication is indeed well-defined. The rest is straightforward.
\end{proof}

\begin{example}[Examples of Quotients]
    \label{ex:quotients}%
    \begin{itemize}
        \item We have ideals \(n \ZZ \idealsub \ZZ\). The quotient ring consists of the cosets
        \[
            \flatfrac{\ZZ}{n \ZZ} = \set{m + n \ZZ}{m \in \ZZ},
        \]
        and therefore we get \(0 + n \ZZ, 1 + n \ZZ, \ldots, \pqty{n - 1} + n\ZZ\). This ring structure takes the structure of arithmetic modulo \(n\).

        \item We take \(\pqty{x} \idealsub \CC\bqty{x}\). The elements of the quotient \(\flatfrac{\CC\bqty{x}}{\pqty{x}}\) can be represented by elements of the form
        \[
            \pqty{a_0 + a_1 x + \cdots + a_n x^n} + \pqty{x}.
        \]

        But the terms \(a_1 x, a_2 x^2, \ldots, a_n x^n \in \pqty{x}\), so this coset is equal to \(a_0 + \pqty{x}\). The representation is unique, so in fact
        \[
            \flatfrac{\CC \bqty{x}}{\pqty{x}} \isomorphic \CC.
        \]
    \end{itemize}
\end{example}

\begin{proposition}[Euclidean Algorithm for Polynomials]
    \label{prop:euclidean-algorithm-polynomials}%
    Let \(K\) be a field. If \(f, g \in K \bqty{x}\), then there exists \(r, q \in K \bqty{x}\), such that
    \[
        f = g \cdot q + r
    \]
    with \(\deg \pqty{r} < \deg \pqty{g}\).
\end{proposition}

\begin{proof}
    Let \(n = \deg \pqty{f}\), so \(f = \sum_{i = 0}^{n} a_i x^i\) where \(a_i \in K\) and \(a_n \neq 0\). Similarly, let \(g = \sum_{i = 0}^{m} b_i x_i\) where \(b_i \in K\) and \(b_m \neq 0\).

    If \(n < m\), set \(q = 0\) and \(r = f\) (purely remainder), and we are done.

    If \(n \geq m\), we proceed by induction on \(n\). Let \(f_1 = f - a_n b_m \inverse x^{n - m} g\). Note in particular, that \(b_m \inverse\) means we are using \(K\) is a \textbf{field}. So
    \[
        \deg \pqty{f_1} < n.
    \]

    If \(n = m\) then \(\deg \pqty{f_1} < n = m\), and hence
    \[
        f = \pqty{a_n b_m \inverse x^{n - m}} g + f_1.
    \]

    If \(n > m\), by induction, we can write \(f_1 = g q_1 + r_1\), where \(\deg \pqty{r_1} < \deg \pqty{g}\). Then,
    \[
        f = a_n b_m \inverse x^{n - m} g + q_1 g + r_1 = \pqty{a_n b_m \inverse x^{n - m} + q_1} g + r_1
    \]
    and we are done.
\end{proof}

A corollary of this is as follows: all ideals in \(K\bqty{x}\) for some field \(K\) are principal. This was shown in \autoref{ex:ideals-of-z}, and we may replace the Euclidean algorithm for \(\ZZ\) with \(K \bqty{x}\). However, note that \(\ZZ \bqty{x}\) and \(K \bqty{x, y}\) (where \(K\) is a field) \textbf{do not} have this property.

Similar to the isomorphism theorems we had in groups, we have them in rings as well.

\begin{theorem}[First Isomorphism Theorem for Rings]
    \label{thm:first-isomorphism-theorem-ring}%
    Let \(\functiondomain{\phi}{R}{S}\) be a ring homomorphism. Then the map
    \[
        \function{f}{\flatfrac{R}{\ker \phi}}{\img \phi \subring S}{r + \ker \phi}{\phi\pqty{r}}
    \]
    is a well-defined ring isomorphism.
\end{theorem}

\begin{proof}
    Bijectivity and additive property follows from \autoref{thm:first-isomorphism-theorem-groups} (First Isomorphism Theorem for Groups). It remains to show this is a ring homomorphism for multiplication. Indeed,
    \[
        f\pqty{\pqty{r + \ker \phi} \cdot \pqty{t + \ker \phi}} = f \pqty{rt + \ker \phi} = \phi \pqty{rt} = \phi\pqty{r} \phi\pqty{t} = f\pqty{r + \ker \phi} \cdot f\pqty{t + \ker \phi}.
    \]
\end{proof}

\begin{example}[Application of First Isomorphism Theorem]
    Consider the homomorphism
    \[
        \function{\phi}{\RR\bqty{x}}{\CC}{f\pqty{x}}{f\pqty{i}}.
    \]

    This map is surjective since \(\phi\pqty{a + bx} = a + bi\). 
    
    The kernel is exactly polynomials such that \(f\pqty{i} = 0\). But recall since this is a polynomial with real coefficients, if \(i\) is a root of \(f\), then so is \(\pqty{-i}\). So \(\pqty{x^2 + 1}\) divides \(f\pqty{x}\). Therefore, the kernel of \(\phi\) is \emph{exactly} \(\pqty{x^2 + 1} \idealsub \RR\bqty{x}\). Therefore, we have
    \[
        \CC \isomorphic \flatfrac{\RR\bqty{x}}{\pqty{x^2 + 1}}.
    \]
\end{example}

\begin{theorem}[Second Isomorphism Theorem for Rings]
    \label{thm:second-isomorphism-theorem-ring}%
    Let \(R \subring S\) and \(J \idealsub S\). Then \(R \cap J \idealsub R\), and \(\flatfrac{\pqty{R + J}}{J} = \set{r + J}{r \in R} \subring \flatfrac{S}{J}\). Furthermore,
    \[
        \flatfrac{R}{R \cap J} \isomorphic \flatfrac{\pqty{R + J}}{J}.
    \]
\end{theorem}

\begin{proof}
    Define a function
    \[
        \function{\phi}{R}{\flatfrac{S}{J}}{r}{r + J}.
    \]

    The kernel of this map is
    \[
        \ker \phi = \set{r \in R}{r + J = 0} = \set{r \in R}{r \in J} = R \cap J.
    \]

    We also have
    \[
        \img \phi = \set{r + J}{r \in R} = \flatfrac{\pqty{R + J}}{J}.
    \]

    By \autoref{thm:first-isomorphism-theorem-ring} (First Isomorphism Theorem for Rings), we have
    \[
        \flatfrac{R}{R \cap J} \isomorphic \flatfrac{\pqty{R + J}}{J}
    \]
    as rings.
\end{proof}

\begin{remark}
    Similar to rings, there is a correspondence for rings and ideals. Let \(I \idealsub R\) be an ideal. There is a bijection
    \[
        \function{f}{\Bqty{\text{subrings of }\flatfrac{R}{I}}}{\Bqty{\text{subrings of }R\text{ containing }I}}{L \subring \flatfrac{R}{I}}{\set{r \in R}{r + I \in L}}
    \]
    and the inverse is given by
    \[
        \function{f\inverse}{\Bqty{\text{subrings of }R\text{ containing }I}}{\Bqty{\text{subrings of }\flatfrac{R}{I}}}{I \idealsub S \subring R}{\flatfrac{S}{I} \subring \flatfrac{R}{I}.}
    \]
\end{remark}

\begin{theorem}[Third Isomorphism Theorem for Rings]
    \label{thm:third-isomorphism-theorem-ring}%
    Let \(I \idealsub R\) and \(J \idealsub R\), with \(I \subseteq J\). Then we have
    \[
        \flatfrac{J}{I} \idealsub \flatfrac{R}{I} \qand \flatfrac{\pqty{\flatfrac{R}{I}}}{\pqty{\flatfrac{J}{I}}} \isomorphic \flatfrac{R}{J}.
    \]
\end{theorem}

\begin{proof}
    Define the function
    \[
        \function{f}{\flatfrac{R}{I}}{\flatfrac{R}{J}}{r + I}{r + J.}
    \]

    It is well-defined by considering this as a group homomorphism.

    This is a ring homomorphism, with kernel
    \[
        \ker f = \set{r + I}{r + J = 0} = \flatfrac{J}{I}.
    \]

    Applying \autoref{thm:first-isomorphism-theorem-ring} (First Isomorphism Theorem for Rings) gives us the result, as desired.
\end{proof}

\begin{proposition}[Ring Homomorphism from \(\ZZ\), Characteristic of a Ring]
    \label{prop:ring-homomorphism-z-characteristic-of-ring}%
    Let \(R\) be any ring. Then there is a homomorphism
    \[
        \functiondomain{i}{\ZZ}{R}
    \]
    sending \(1\) to \(1_R\). For \(i\), we have
    \[
        \ker\pqty{i} = n \ZZ
    \]
    for some \(n \in \ZZ\), where \(\abs{n}\) is known as the \emph{characteristic} of \(R\).
\end{proposition}

Therefore, \(\ZZ\) is sometimes called the \emph{initial ring}. Furthermore, note that the zero ring does \textbf{not} have any homomorphisms to another ring (not a single one) -- this is different role from the trivial group back for group homomorphism.

\begin{example}[Characteristics of Common Rings]
    The infinite number rings \(\ZZ, \RR, \CC\) all have characteristic \(0\). The ring \(\flatfrac{\ZZ}{k\ZZ}\) has characteristic \(k\).
\end{example}